\documentclass[10pt,a4paper]{amsart}
\usepackage[margin=19mm]{geometry}
\usepackage{amsmath,amssymb,mathtools}
\usepackage[scr=rsfs,cal=euler]{mathalfa}
\usepackage{xfrac}
\usepackage[unicode,hidelinks,bookmarks=false]{hyperref}
\newcommand{\ie}{i.e.\ }
\newcommand{\cf}{cf.\ }
\newcommand{\resp}{respectively\ }
\newcommand{\Subsection}{\subsection}
\providecommand{\nobreakdash}{\nobreak\hbox{-}}
\setlength{\emergencystretch}{2em}
\multlinegap0pt
\usepackage{amssymb}
\newtheorem{theorem}{Theorem}
\title{Lyra Almost Ricci-Bourguignon Solitons on Twisted Warped Product Manifolds}
\author{Ayman Elsharkawy, Eman Ghareeb Rezk, Uday Chand De, Emad F. Wanas, Abdelrahman M. Tawfiq}
\date{Source: arXiv:2608.14248v1 (CC BY 4.0)}
\begin{document}
\maketitle

\noindent\emph{Source and licence.} Lyra Almost Ricci-Bourguignon Solitons on Twisted Warped Product Manifolds. Version arXiv:2608.14248v1 (CC BY 4.0), distributed by arXiv under the Creative Commons Attribution 4.0 International licence, http://creativecommons.org/licenses/by/4.0/, as recorded in the arXiv licence metadata for this exact version and shown on the abstract page https://arxiv.org/abs/2608.14248v1. Full licence notice: \texttt{licenses/arxiv-2608.14248-CC-BY-4.0.txt}.\par\smallskip

\noindent\textit{Editorial scope.} Counted unit: the article's theorem thm:conformal-einstein-equivalence (source0.tex lines 3164--3209). The article's complete original proof of it is delivered with it (source0.tex lines 3211--3309, one \texttt{proof} environment of 99 lines). No mathematics is written, repaired or re-derived: the statement and the proof are the article's own lines, in the article's own notation, and no retained line is substituted or re-flowed.  Retained uncounted as the source-local context the proof needs: lines 838--847.  Nothing was omitted from any retained span, so the package carries no omission record.  No retained line contains a citation, so the wrapper carries no bibliography and no bibliography entry.  No retained line contains a figure environment, an image-inclusion command or drawing code, so no image is delivered and the package declares no asset dependency.  Editorial formatting only, in addition: an AMS \texttt{amsart} preamble that restates the article's own declarations and macros used by the retained lines, namely the environments \texttt{theorem} on separate counters, together with the standard packages those lines need.  The retained environments print in this document as Theorem 1 in order of appearance, whereas the article prints the same items with the numbering of its own full document.  The complete native source of the article as distributed in the arXiv e-print for this exact version is bundled unchanged as \texttt{sources/207652/source0.tex}.
	\begin{equation}
		\operatorname{Ric}^{L}
		+
		\frac{1}{2}\mathcal L_{Z}g_{L}
		=
		\left(
		\lambda+\rho R^{L}
		\right)g_{L},
		\label{eq:lyra-arbs-equation}
	\end{equation}

\begin{theorem}[Conformal--Einstein equivalence]
	\label{thm:conformal-einstein-equivalence}
	Let
	\[
	\left(
	M,g_{L},Z,\lambda,\rho
	\right)
	\]
	be a Lyra almost Ricci--Bourguignon soliton on a connected
	manifold of dimension \(n\geq3\). Then the following statements
	are equivalent:
	\begin{enumerate}
		\item \(Z\) is Lyra conformal;
		\item \(g_{L}\) is an Einstein metric.
	\end{enumerate}
	
	More precisely, if
	\[
	\mathcal L_{Z}g_{L}=2\psi g_{L},
	\]
	then
	\begin{equation}
		\operatorname{Ric}^{L}
		=
		\mu g_{L},
		\qquad
		\mu
		=
		\lambda+\rho R^{L}-\psi,
		\label{eq:einstein-factor-conformal}
	\end{equation}
	where \(\mu\) is constant. Moreover,
	\begin{equation}
		R^{L}
		=
		n\mu
		\label{eq:einstein-scalar-curvature}
	\end{equation}
	and
	\begin{equation}
		\lambda
		=
		(1-n\rho)\mu+\psi.
		\label{eq:lambda-conformal-einstein}
	\end{equation}
\end{theorem}

\begin{proof}
	Suppose that \(Z\) is Lyra conformal. Substituting
	\[
	\frac12\mathcal L_{Z}g_{L}
	=
	\psi g_{L}
	\]
	into \eqref{eq:lyra-arbs-equation} gives
	\[
	\operatorname{Ric}^{L}
	=
	\left(
	\lambda+\rho R^{L}-\psi
	\right)g_{L}.
	\]
	Define
	\[
	\mu
	=
	\lambda+\rho R^{L}-\psi.
	\]
	Then
	\[
	\operatorname{Ric}^{L}
	=
	\mu g_{L}.
	\]
	Taking the trace yields
	\[
	R^{L}=n\mu.
	\]
	
	The contracted second Bianchi identity gives
	\[
	\operatorname{div}_{L}
	\operatorname{Ric}^{L}
	=
	\frac12 dR^{L}.
	\]
	Since
	\[
	\operatorname{div}_{L}(\mu g_{L})=d\mu
	\]
	and
	\[
	dR^{L}=n\,d\mu,
	\]
	we obtain
	\[
	d\mu
	=
	\frac n2 d\mu.
	\]
	Because \(n\geq3\),
	\[
	\left(
	1-\frac n2
	\right)d\mu=0
	\]
	implies
	\[
	d\mu=0.
	\]
	Thus, \(\mu\) is constant and \(g_{L}\) is Einstein.
	
	Conversely, suppose that
	\[
	\operatorname{Ric}^{L}
	=
	\mu g_{L}.
	\]
	Using \eqref{eq:lyra-arbs-equation}, we find
	\[
	\frac12\mathcal L_{Z}g_{L}
	=
	\left(
	\lambda+\rho R^{L}-\mu
	\right)g_{L}.
	\]
	Hence \(Z\) is Lyra conformal with factor
	\[
	\psi
	=
	\lambda+\rho R^{L}-\mu.
	\]
	
	Finally, substituting
	\[
	R^{L}=n\mu
	\]
	into \eqref{eq:einstein-factor-conformal} gives
	\[
	\mu
	=
	\lambda+n\rho\mu-\psi,
	\]
	which is equivalent to
	\eqref{eq:lambda-conformal-einstein}.
\end{proof}

\end{document}
